\section{模型假设}
三个问题使用同一光学评价器，问题二和问题三只改变设计变量与约束，并不重新定义太阳辐射、反射和能量聚合方式。因此，下列假设同时适用于三问；坐标范围、给定镜场规模和集热器尺寸等题面事实不列作假设。

\begin{enumerate}[itemindent=2em]
  \item \textbf{全年晴空假设。}假设全年规定评价时段均为晴天，不考虑云层对太阳直射辐射的随机遮蔽，法向直接辐射强度统一由题目给定的DNI公式计算。该假设使各月规定时刻能够在相同气象口径下比较，但所得年平均功率不包含阴雨天气和云量变化造成的实际可利用率损失。
  \item \textbf{理想平面镜假设。}定日镜没有题面未给出的面形、斜率和跟踪误差，镜面法向使太阳中心光线指向集热器中心。该假设使反射几何可唯一确定，但真实电站的截断效率可能更低。
  \item \textbf{均匀太阳盘假设。}太阳锥采用半角$4.65\times10^{-3}$ rad、单位立体角能量均匀的pillbox分布。题面没有提供周日比，故不额外引入需要经验参数的周日区模型；问题一再用半角$\pm5\%$检验该简化对功率的影响。
  \item \textbf{联合光线计损假设。}阴影和遮挡均按每条光线的布尔可见性判断，多镜相交不重复计损；截断效率以未受阴影遮挡的反射光为条件分母。这是使题面效率乘积符合能量守恒的必要口径。
  \item \textbf{离散年平均假设。}每月五个规定时点等权形成月平均，十二个月等权形成题面离散年平均。三问均采用这一口径，使评价与优化目标一致；它不等同于逐日全年积分，问题一另与按月天数加权结果比较。
  \item \textbf{塔身与集热器几何。}三问均把塔身表示为半径3.5 m、$0\le z<72$ m的实体圆柱并计入入射阴影；集热器侧面取$72\le z\le80$ m，只在反射后的命中判定中使用。到达集热器侧面的反射能量按题面公式计入输出热功率，不额外加入未给出的吸收率或热损失。
\end{enumerate}
